\documentclass[11pt,fontset=fandol]{ctexart}
\usepackage[a4paper,margin=23mm]{geometry}
\usepackage{amsmath,amssymb,mathtools,bm,mathrsfs,esint}
\usepackage{graphicx,longtable,booktabs,array,calc}
\usepackage[normalem]{ulem}
\usepackage{xcolor,hyperref}
\hypersetup{hidelinks}
\setcounter{secnumdepth}{-1}
\newcounter{none}
\providecommand{\tightlist}{\setlength{\itemsep}{0pt}\setlength{\parskip}{0pt}}
\providecommand{\passthrough}[1]{#1}
\usepackage{amssymb,mathrsfs,esint,bm,newunicodechar}
\allowdisplaybreaks
\setlength{\emergencystretch}{3em}
\providecommand{\llangle}{\langle\!\langle}
\providecommand{\rrangle}{\rangle\!\rangle}
\newunicodechar{∞}{\ensuremath{\infty}}
\newunicodechar{∈}{\ensuremath{\in}}
\newunicodechar{∑}{\ensuremath{\sum}}
\newunicodechar{∫}{\ensuremath{\int}}
\newunicodechar{≤}{\ensuremath{\leq}}
\newunicodechar{≥}{\ensuremath{\geq}}
\newunicodechar{≠}{\ensuremath{\neq}}
\newunicodechar{→}{\ensuremath{\to}}
\newunicodechar{⇒}{\ensuremath{\Rightarrow}}
\newunicodechar{α}{\ensuremath{\alpha}}
\newunicodechar{β}{\ensuremath{\beta}}
\newunicodechar{γ}{\ensuremath{\gamma}}
\newunicodechar{λ}{\ensuremath{\lambda}}
\newunicodechar{φ}{\ensuremath{\varphi}}
\newunicodechar{ε}{\ensuremath{\varepsilon}}
\newunicodechar{σ}{\ensuremath{\sigma}}
\newunicodechar{θ}{\ensuremath{\theta}}
\newunicodechar{δ}{\ensuremath{\delta}}
\newunicodechar{η}{\ensuremath{\eta}}
\newunicodechar{π}{\ensuremath{\pi}}
\newunicodechar{ξ}{\ensuremath{\xi}}
\newunicodechar{Σ}{\ensuremath{\Sigma}}
\newunicodechar{Δ}{\ensuremath{\Delta}}
\newunicodechar{∂}{\ensuremath{\partial}}
\newunicodechar{∀}{\ensuremath{\forall}}
\newunicodechar{∃}{\ensuremath{\exists}}
\newunicodechar{√}{\ensuremath{\surd}}
\newunicodechar{−}{\ensuremath{-}}
\newunicodechar{μ}{\ensuremath{\mu}}
\newunicodechar{ρ}{\ensuremath{\rho}}
\newunicodechar{τ}{\ensuremath{\tau}}
\newunicodechar{ω}{\ensuremath{\omega}}
\newunicodechar{Ω}{\ensuremath{\Omega}}
\newunicodechar{χ}{\ensuremath{\chi}}
\newunicodechar{ψ}{\ensuremath{\psi}}
\newunicodechar{Γ}{\ensuremath{\Gamma}}
\newunicodechar{Φ}{\ensuremath{\Phi}}
\newunicodechar{Λ}{\ensuremath{\Lambda}}
\newunicodechar{ℝ}{\ensuremath{\mathbb{R}}}
\newunicodechar{ℤ}{\ensuremath{\mathbb{Z}}}
\newunicodechar{ℕ}{\ensuremath{\mathbb{N}}}
\newunicodechar{△}{\ensuremath{\triangle}}
\newunicodechar{⊗}{\ensuremath{\otimes}}
\begin{document}
\section{2016年杭州电子科技大学601数学分析考研真题}\label{ux5e74ux676dux5ddeux7535ux5b50ux79d1ux6280ux5927ux5b66601ux6570ux5b66ux5206ux6790ux8003ux7814ux771fux9898}

来源：\texttt{HDU/601数学分析/2016/2016年杭州电子科技大学601数学分析考研真题.pdf}。原
PDF 页序：1, 2。

\section{杭州电子科技大学2016年攻读硕士学位研究生招生考试}\label{ux676dux5ddeux7535ux5b50ux79d1ux6280ux5927ux5b662016ux5e74ux653bux8bfbux7855ux58ebux5b66ux4f4dux7814ux7a76ux751fux62dbux751fux8003ux8bd5}

\section{《数学分析》试题}\label{ux6570ux5b66ux5206ux6790ux8bd5ux9898}

（试题共四大题，共2页，总分150分）

姓名 \_\_\_\_　报考专业 \_\_\_\_　准考证号 \_\_\_\_

【所有答案必须写在答题纸上，做在试卷或草稿纸上无效！】

\begin{quote}
助手整理依据：交叉核对远程OCR、本地逐页OCR和同题原参考答案，恢复丢失的公式语法。二4终点、四1不等式均保留原印条件；三2积分区域由本地OCR恢复，原极限变量与积分参数冲突的情况明确标出。
\end{quote}

\subsection{一、填空题（8小题，每题5分，共40分）}\label{ux4e00ux586bux7a7aux98988ux5c0fux9898ux6bcfux98985ux5206ux517140ux5206}

\begin{enumerate}
\def\labelenumi{\arabic{enumi}.}
\item
  求 \(\displaystyle\lim_{x\to0}\frac{\sqrt{1+\sin^2x}-1}{x\ln(1+x)}=\)
  \_\_\_\_。
\item
  设 \[
  f(x)=\begin{cases}x^\alpha\sin(1/x),&x\ne0,\\0,&x=0.\end{cases}
  \] 则实数 \(\alpha\) \_\_\_\_ 时 \(f\) 在零点连续，\(\alpha\) \_\_\_\_
  时 \(f\) 在零点可导（写出范围）。
\item
  若 \(\displaystyle f(x)=\frac d{dx}\int_0^x\cos(t+x)\,dt\)，则
  \(f(x)=\) \_\_\_\_。
\item
  \(\displaystyle\int_0^{80\pi}\sqrt{1-\cos2x}\,dx=\) \_\_\_\_。
\item
  设 \(y(x)=(x-2)/(3x+1)\)，则 \(y^{(n)}(x)=\) \_\_\_\_。
\item
  设 \(f(x,y)\) 连续，交换
  \(\displaystyle\int_0^a dx\int_0^xf(x,y)\,dy\) 的积分次序，得到
  \_\_\_\_。
\item
  级数 \(\displaystyle\sum_{n=1}^\infty\frac{(x-2)^{2n}}{n4^n}\)
  的收敛域为 \_\_\_\_。
\item
  设 \(y=y(x)\) 由 \(e^{x+y}+\cos(xy)=0\) 确定，则 \(dy/dx=\) \_\_\_\_。
\end{enumerate}

\subsection{二、计算题（4小题，共30分）}\label{ux4e8cux8ba1ux7b97ux98984ux5c0fux9898ux517130ux5206}

\begin{enumerate}
\def\labelenumi{\arabic{enumi}.}
\item
  （5分）求 \(\displaystyle\lim_{x\to0}\frac{3^x+3^{-x}-2}{x^2}\)。
\item
  （5分）求 \(\displaystyle\int e^{\sin x}\sin2x\,dx\)。
\item
  （10分）将 \(f(x)=(1+x)/(1-x)^3\) 在 \(x=0\) 展成幂级数，并求
  \(\displaystyle\sum_{n=1}^\infty n^2/2^{n-1}\)。
\item
  （10分）计算 \[
  I=\int_L(5xy-e^x\sin y)\,dy+e^x\cos y\,dx,
  \] 其中 \(L\) 为曲线 \(x=\sqrt{2y-y^2}\) 从 \(A(0,0)\) 到 \(B(2,0)\)
  的一段。

  原印终点 \(B(2,0)\) 不在所给曲线上；原参考答案改为 \(B(0,2)\)
  后计算。此处保留原题，修正版的计算单独见解析。
\end{enumerate}

\subsection{三、综合题（3小题，每题10分，共30分）}\label{ux4e09ux7efcux5408ux98983ux5c0fux9898ux6bcfux989810ux5206ux517130ux5206}

\begin{enumerate}
\def\labelenumi{\arabic{enumi}.}
\item
  判别
  \(\displaystyle\sum_{n=1}^\infty\frac{(-1)^{n-1}\sin(nx)}{\sqrt[3]{n^5+x^4}}\)
  在 \(\mathbb R\)
  上的一致收敛性，说明理由。分母按远程题面及原答案共同识别为三次根式；本地OCR的根号范围识读有误。
\item
  设 \(f(u)\) 在零点某邻域内连续，在零点可导且
  \(f(0)=0\)。讨论以下极限是否存在，若存在则求值： \[
  \lim_{t\to0}\frac1{\pi t^4}\iiint_{x^2+y^2+z^2\le t^2}
  f(\sqrt{x^2+y^2+z^2})\,dv.
  \] 原题两份OCR均将极限下标识作 \(x\to0\)，而 \(t\)
  才是球域半径参数、\(x\)
  已是积分哑变量；此处依据原参考答案的完整写法整理为 \(t\to0\)。
\item
  讨论
  \(\displaystyle\int_1^\infty\frac{\sin x}{x^p}\,dx\)（\(p>0\)）绝对收敛和条件收敛时
  \(p\) 的范围。
\end{enumerate}

\subsection{四、证明题（5小题，每题10分，共50分）}\label{ux56dbux8bc1ux660eux98985ux5c0fux9898ux6bcfux989810ux5206ux517150ux5206}

\begin{enumerate}
\def\labelenumi{\arabic{enumi}.}
\item
  证明当 \(x>0\) 时 \(\displaystyle\frac1{1+x}<\ln(1+x)<x\)。

  此不等式按原印保留，原参考答案指出左端有误，并证明
  \(x/(1+x)<\ln(1+x)<x\)。
\item
  设 \(f\) 可导，证明曲面 \(z=xf(y/x)\) 的切平面通过原点；原题另规定
  \(x=0\) 时 \(z=0\)。
\item
  设 \(f\) 在 \([0,1]\) 连续、\((0,1)\) 可导，\(f(0)=0\)。证明存在
  \(\xi\in(0,1)\)，使 \(\xi f'(\xi)+2f(\xi)=f'(\xi)\)。
\item
  证明 \(f(x)=\sin(1/x)\) 在 \((0,1)\) 连续但不一致连续。
\item
  设 \(f\) 在 \([0,1]\) 严格单调递减且连续。证明任意 \(q\in(0,1)\) 均有
  \[
  \int_0^qf(x)\,dx>q\int_0^1f(x)\,dx.
  \]
\end{enumerate}
\end{document}
